Write \(h=N^{-1}\) and choose a fixed compact annulus
\(c_0\le|\eta|\le C_0\) containing the support of \(\varphi\).
Choose a smooth finite partition \(\varphi=\sum_\nu\varphi_\nu\)
subordinate to caps of the aperture in \ref{ms:prop-6-3} for this
annulus. On the torus \(\mathbb R^n/(100\mathbb Z)^n\), apply
\ref{ms:prop-6-3} to \(\varphi_\nu(hD)F\). Its \(L^{p_c}\) norm is
bounded uniformly by \(C_\nu\|F\|_{p_c}\): its convolution kernel
is the periodization of \(h^{-n}\check\varphi_\nu(\cdot/h)\),
whose \(L^1\) norm is at most \(\|\check\varphi_\nu\|_1\).
Here \(\check\varphi_\nu\) denotes the inverse Fourier transform.
The torus carries ordinary Lebesgue volume on a fundamental cube,
so the periodized kernel has the same multiplier at frequencies
\((2\pi/100)\mathbb Z^n\). Summing over \(\nu\) yields
\[
 \|U(T)\varphi(hD)F\|_{L^{p_c}(I\times\mathbb T^n)}
 \le C[\log(2/h)]^{2/p_c}\|F\|_{L^{p_c}(\mathbb T^n)}.
\]

On \(\mathbb R^n\) the convolution kernel is
\[
 K_h(T,X)=(2\pi h)^{-n}
   \int e^{i(X\cdot\eta-T|\eta|)/h}\varphi(\eta)\,d\eta .
\]
For \(T\in I\) and \(|X|\ge3\), the phase \(\Phi=X\cdot\eta-T|\eta|\)
satisfies \(|\nabla_\eta\Phi|\ge |X|/3\), and its derivatives of
order at least two are uniformly bounded on the annulus. Integrate
by parts with \(h\nabla_\eta\Phi\cdot\nabla_\eta/(i|\nabla_\eta\Phi|^2)\).
The coefficient without \(h\) and all of its derivatives on the
support are \(O(|X|^{-1})\). Thus \(M+n\) iterations for integer
\(M\ge1\) give
\[
 |K_h(T,X)|\le C_M h^M(1+|X|)^{-M}
 \qquad (T\in I,\ |X|\ge3).
\]
Tile \(\mathbb R^n\) by unit cubes \(Q\), and let \(10Q\) be the
concentric cube of side ten. For \(X\in Q\) and \(Y\notin 10Q\)
one has \(|X-Y|\ge9/2\). On every cube the far contribution is
therefore bounded by the single function
\[
 C_M h^M\big((1+|\cdot|)^{-M}*|f|\big)(X).
\]
For \(M>n\), this function has \(L^{p_c}(I\times\mathbb R^n)\)
norm at most \(C_M h^M\|f\|_{p_c}\) by Young's inequality and
\(|I|=1\).

For the near contribution put \(f_Q=f\mathbf 1_{10Q}\). Translate the
center of \(Q\) to the origin and periodize \(f_Q\) at period \(100\).
The periodic input has norm \(\|f\|_{L^{p_c}(10Q)}\). The periodic
output uses \(\sum_{k\in\mathbb Z^n}K_h(T,X+100k)\) as its kernel.
For \(X\in Q\), \(Y\in 10Q\), and \(k\ne0\), the argument
\(X-Y+100k\) has length at least \(94.5|k|\): for every nonzero
coordinate \(k_j\), \(|X_j-Y_j+100k_j|\ge 94.5 |k_j|\).
The decay bound and \(M>n\) show that the periodic and Euclidean
outputs differ on \(I\times Q\) by at most
\(C_M h^M\|f_Q\|_1\) pointwise. The volume of \(10Q\) is fixed, so
the torus bound and Hölder's inequality give
\[
 \|U(T)P_Nf_Q\|_{L^{p_c}(I\times Q)}
 \le C[\log(2/h)]^{2/p_c}\|f\|_{L^{p_c}(10Q)}.
\]
Take \(p_c\)-th powers and sum over \(Q\). The cubes \(10Q\) have
bounded overlap. Adding the far output proves the result since
\(\log(2/h)\) is comparable to \(\log(2+N)\).
